\documentclass[11pt]{article}
\usepackage[margin=1in]{geometry}
\usepackage{amsmath}
\usepackage{amssymb}
\usepackage{enumitem}
\usepackage{xcolor}
\usepackage{tikz}
\usepackage{booktabs}
\usepackage{tcolorbox}

% Custom colors
\definecolor{boxblue}{RGB}{230,240,255}
\definecolor{boxgreen}{RGB}{230,255,240}

\title{\textbf{Problem Set 3, Question 2(a)} \\ Short-Run Equilibrium Explained}
\author{ECO310: Math Intermediate Microeconomics}
\date{Spring 2026}

\begin{document}
\maketitle

\section{What is Short-Run Equilibrium?}
\textit{(Lecture 11-12, slide 31)}

\begin{tcolorbox}[colback=boxblue, colframe=blue!50!black, title=From the slides:]
``In equilibrium, demand equals supply. Equilibrium price: the price $p$ at which $Q^S(p) = Q^D(p)$''

\vspace{0.3cm}
``This is the \textbf{ONLY} requirement for a SR equilibrium''
\end{tcolorbox}

\subsection*{In Plain English:}

\textbf{Short-run equilibrium} is when the market clears:
\begin{itemize}[leftmargin=*]
    \item \textbf{Supply = Demand}
    \item The number of firms $N$ is \textbf{fixed} (firms can't enter or exit in the short run)
    \item Each firm produces whatever maximizes their profit at the market price
\end{itemize}

\section{Problem 2(a) Setup}

\subsection*{Given:}
\begin{itemize}[leftmargin=*]
    \item Market demand: $Q^D(p) = 200 - 5p$
    \item $N = 10$ firms (fixed in short run)
    \item Each firm has cost: $C(q) = f(N) + \alpha q^2$
    \begin{itemize}
        \item For part (a): we're just told $N = 10$, and $\alpha > 0$
        \item $f(N)$ is a ``fixed cost'' that depends on how many firms are in the market
    \end{itemize}
\end{itemize}

\subsection*{Find:} The short-run equilibrium price

\newpage

\section{Solution Steps}
\textit{(Following Lecture 11-12, slides 17-21, 29-31)}

\subsection{Step 1: Find each firm's supply curve}

Each \textbf{price-taking firm} maximizes profit by producing where \textbf{$p = MC$} (Lecture 11-12, slide 7)

\vspace{0.3cm}
\textbf{Cost function:}
\[
C(q) = f(N) + \alpha q^2
\]

\textbf{Marginal cost:}
\[
MC(q) = \frac{dC}{dq} = 2\alpha q
\]

\textbf{Each firm's supply curve} (Lecture 11-12, slide 20):
\begin{align*}
p &= MC \\
p &= 2\alpha q \\
q &= \frac{p}{2\alpha} \quad \leftarrow \text{ This is each firm's supply curve } q^S(p)
\end{align*}

\subsection{Step 2: Find market supply}
\textit{(Lecture 11-12, slides 29-30)}

\begin{tcolorbox}[colback=boxgreen, colframe=green!50!black]
\textbf{From the slides:} ``To derive market supply, we add up the individual market supply curves across firms''
\end{tcolorbox}

With $N = 10$ firms:
\[
Q^S(p) = N \times q^S(p) = 10 \times \frac{p}{2\alpha} = \frac{5p}{\alpha}
\]

\subsection{Step 3: Set supply = demand}
\textit{(Lecture 11-12, slide 31)}

\textbf{Equilibrium condition:}
\begin{align*}
Q^S(p) &= Q^D(p) \\
\frac{5p}{\alpha} &= 200 - 5p \\
\frac{5p}{\alpha} + 5p &= 200 \\
5p\left(\frac{1}{\alpha} + 1\right) &= 200 \\
5p\left(\frac{1 + \alpha}{\alpha}\right) &= 200 \\
p &= \frac{200\alpha}{5(1 + \alpha)} \\
p^* &= \frac{40\alpha}{1 + \alpha}
\end{align*}

\begin{tcolorbox}[colback=yellow!10, colframe=orange!75!black, title=Final Answer:]
\Large
\[
\boxed{p^* = \frac{40\alpha}{1 + \alpha}}
\]
\end{tcolorbox}

\newpage

\section{Key Differences: Short Run vs. Long Run}

\begin{center}
\begin{tabular}{p{6cm}|p{6cm}}
\toprule
\textbf{Short Run (SR)} & \textbf{Long Run (LR)} \\
\midrule
Number of firms $N$ is \textbf{fixed} & Firms can \textbf{enter/exit} \\
\addlinespace
Only condition: Supply = Demand & Additional condition: Firms make \textbf{zero profit} \\
\addlinespace
Firms might make profit or loss & $p = \min AC$ (firms break even) \\
\bottomrule
\end{tabular}
\end{center}

\vspace{0.5cm}

\begin{tcolorbox}[colback=boxblue, colframe=blue!50!black]
\textbf{From slide 32:} ``If firms currently in the market are making money, new firms will come in... if firms are losing money, some will exit''
\end{tcolorbox}

\vspace{0.3cm}

In the \textbf{short run}, we don't worry about profits yet --- we just find where supply meets demand with the current number of firms!

\section{Intuition Check}

Think of it like a farmers market:

\begin{itemize}[leftmargin=*]
    \item \textbf{Short run:} There are 10 farmers at the market today (can't change)
    \item Each farmer decides how much to produce based on the market price
    \item The price adjusts until the total amount farmers bring = total amount customers want
    \item Some farmers might make money, some might lose money, but they're stuck there for today
\end{itemize}

\vspace{0.5cm}

\textit{That's short-run equilibrium!}

\end{document}
